\documentclass[11pt,letterpaper]{article}
\usepackage[margin=1.02in]{geometry}
\usepackage[T1]{fontenc}
\usepackage{lmodern}
\usepackage{microtype}
\usepackage{amsmath,amssymb,amsthm,mathtools,mathrsfs}
\usepackage{booktabs,array,enumitem,xcolor,seqsplit}
\usepackage[colorlinks=true,linkcolor=blue!45!black,citecolor=blue!40!black,urlcolor=blue!40!black]{hyperref}
\usepackage{fancyhdr}
\pagestyle{fancy}\fancyhf{}\lhead{NIGHTFALL \quad v9.2}\rhead{Beyond the Boolean cube}\cfoot{\thepage}
\setlength{\headheight}{14pt}
\setlist{itemsep=2pt,topsep=3pt}
\newtheorem{theorem}{Theorem}[section]
\newtheorem{proposition}[theorem]{Proposition}
\newtheorem{lemma}[theorem]{Lemma}
\newtheorem{corollary}[theorem]{Corollary}
\newtheorem{remark}[theorem]{Remark}
\newtheorem{question}[theorem]{Question}
\theoremstyle{definition}\newtheorem{definition}[theorem]{Definition}
\newcommand{\E}{\mathbb E}
\newcommand{\Var}{\operatorname{Var}}
\newcommand{\osc}{\operatorname{osc}}
\newcommand{\Inf}{\operatorname{I}}
\newcommand{\cH}{\mathcal H}
\newcommand{\cK}{\mathcal K}
\newcommand{\Ctwo}{C_2}
\newcommand{\cprod}{c_{\rm prod}}
\newcommand{\Law}{\operatorname{Law}}
\newcommand{\supp}{\operatorname{supp}}
\newcommand{\rank}{\operatorname{rank}}
\newcommand{\muprod}{\mu}
\newcommand{\Tnoise}{T_t}
\newcommand{\abs}[1]{\lvert #1\rvert}
\newcommand{\norm}[1]{\lVert #1\rVert}
\newcommand{\status}[1]{\textsf{#1}}
\begin{document}
\begin{center}
{\Large\bfseries Beyond the Boolean Cube:\par}
\vspace{3pt}
{\Large\bfseries Degree-Two Influence and Rank-Optimal Observation\par}
\vspace{3pt}
{\large Bias-Free Theorems on Arbitrary Product Spaces\par}
\vspace{3pt}
{\normalsize With an All-Degree Extension and Scoped Lean Evidence}\par
\vspace{10pt}
Devin J. Wright\\NIGHTFALL research program\\
12 November 2026 \quad\textbf{Version 9.2: degree-two results lead}
\end{center}
\begin{abstract}
\noindent\textbf{Principal result: the degree-two law.}
For every bounded real observable $|f|\le B$ of Hoeffding interaction degree at most two on an arbitrary independent product law, including heterogeneous, arbitrarily biased, categorical and nonatomic inputs, we give a direct matching--Hilbert-fourth-moment proof of
\[
\max_i\Inf_i(f)\ \ge\ \frac{\Var(f)^2}{\Ctwo B^2},
\qquad \Ctwo=(10+2\sqrt2)^2(1+\sqrt3)^2<1229.
\]
The variance exponent two is optimal even for fair binary sources. For rank-$r$ Hilbert-valued quadratic outputs we prove the companion $V(G)^2\le 3\Ctwo rB^2\max_i\Inf_i(G)$; an exact posterior energy identity yields a deterministic adaptive whole-symbol prediction policy with risk at most $\tau_r+\eta$ after at most $\lceil12\Ctwo rB^2/\eta\rceil$ queries (or all coordinates), with matching $\Omega(r/\eta)$ scaling. These are analytic results with written proofs, not end-to-end Lean certifications or polynomial-time policy constructions.

\noindent\textbf{Complementary general-degree extension.}
Separately, for interaction degree $d$, we retain the bias-free product influence estimate $\max_i\Inf_i(f)\ge\cprod\Var(f)^{5/2}/(B^3d^6)$ with $\cprod=0.001201250694\ldots$ from the unrefereed v8 analytic chain. We distinguish this from the weaker constant $c_* = \sqrt{3/12500}/\pi^6$ supported by a pinned finite-product Lean declaration and a separately written finite-to-measurable lift. The Boolean-cube solution of Turturean and the inherited weighted machinery are explicitly credited. We prioritize the degree-two proofs and their recovery consequences, then present the general-degree extension and all certification boundaries. Global priority and independent specialist review remain outstanding.
\end{abstract}
\noindent\textbf{Keywords:} bias-free product-space influence; Hoeffding degree; Efron--Stein decomposition; arbitrary independent measures; quadratic interactions; adaptive observation; Lean verification; Hilbert-valued functions.\\
\textbf{Status:} Specialist-referee draft. No external submission or claim of globally established priority. The degree-two analytic proofs are written in full below; the all-degree result imports named finite and analytic premises, with evidence boundaries in Section~\ref{sec:status}.

\section{Degree-two theorems: influence, optimal exponent, and rank-optimal recovery}

\noindent\textbf{Degree-two first.} The flagship assertions are Theorem~\ref{thm:q} (law-uniform variance-squared influence), Proposition~\ref{prop:sharp} (sharp exponent), Theorem~\ref{thm:r} (rank-linear Hilbert transfer), and Theorem~\ref{thm:obs} (matching-order adaptive observation). They require only the independent-product assumptions stated below, not the all-degree analytic strengthening. The comparison with existing Boolean reductions is deferred until after their proofs.\par\medskip
Turturean \cite{Turturean2026} proved a universal $V^{5/2}/d^6$ influence lower bound on the \emph{uniform Boolean cube}; this paper does not claim a second proof of that result. Girish--Mittal--Nehoran--Raz \cite{GMNR2026} proved equivalences among fixed-bias Boolean product and slice formulations, allowing bias-dependent quantitative transfer. Neither paper, in its stated general theorems, gives the arbitrary-alphabet quadratic and Hilbert-output guarantees below. Earlier work on completely bounded polynomials has a distinct stronger norm hypothesis \cite{EG2023,EGSP2026}; binary quadratic cases also have stronger classical Grothendieck estimates. Thus our emphasis is \emph{the absence of alphabet, atom-mass, and dimension parameters under only pointwise boundedness}, not the originality of the variance exponent by itself.

For finite independent $X=(X_1,\ldots,X_n)$, with laws $\mu_i$ on spaces $\Omega_i$, let $E_i$ integrate coordinate $i$ and $D_i=\mathrm{Id}-E_i$. Set
\begin{equation}\label{eq:infdef}
\Inf_i(f)=\E\Var(f\mid X_{-i})=\norm{D_i f}_{L^2(\mu)}^2
=\tfrac12\E\bigl(f(X)-f(X^{(i)})\bigr)^2,
\end{equation}
where $X^{(i)}$ resamples $i$ independently. For $\cH$-valued functions, replace the square by the squared Hilbert norm. The Hoeffding expansion is $f=\sum_{S\subset[n]}f_S$ with $E_i f_S=0$ for $i\in S$; its interaction degree is at most $d$ if $f_S=0$ for $\abs S>d$. For finite products,
\begin{equation}\label{eq:anova}
\Var(f)=\sum_{S\ne\varnothing}\norm{f_S}_2^2,\qquad
\Inf_i(f)=\sum_{S\ni i}\norm{f_S}_2^2,\qquad
\Var(f)\le\sum_i\Inf_i(f)\le d\Var(f).
\end{equation}
All effective finite alphabets are understood after zero-mass symbols are removed. We use \emph{whole-symbol} observations: one query reveals the entire value of one coordinate, whatever its alphabet.

Define the constants
\begin{equation}\label{eq:constants}
D_2=(10+2\sqrt2)(1+\sqrt3),\qquad
\Ctwo=D_2^2=432+160\sqrt2+216\sqrt3+80\sqrt6
=1228.3563238\ldots.
\end{equation}

\begin{theorem}[Law-uniform quadratic influence]\label{thm:q}
Let $\{X_i\}_{i\in J}$ be an arbitrary independent family, on arbitrary probability spaces, with $J$ of any cardinality. Suppose $f\in L^2(\mu)$ is real, $\abs f\le B$ almost surely for $B>0$, and has Hoeffding interaction degree at most two. If $V=\Var(f)>0$, then the maximum resampling influence is attained and
\begin{equation}\label{eq:Q}
\boxed{\displaystyle \max_{i\in J}\Inf_i(f)\ge\frac{V^2}{\Ctwo B^2}.}
\end{equation}
The constant is independent of biases, atom masses, alphabet sizes, and the number of coordinates.
\end{theorem}

\begin{proposition}[Optimal variance exponent]\label{prop:sharp}
For every $m\ge1$ there exists $f_m:\{-1,1\}^{2m}\to[-1,1]$ of \emph{exact} degree two for which $\Var(f_m)=1/m$ and $\max_i\Inf_i(f_m)=1/m^2$. Thus no uniform positive lower bound $c\Var(f)^\alpha$ can hold for all bounded degree-two product observables when $\alpha<2$.
\end{proposition}

\begin{theorem}[Rank-linear quadratic influence]\label{thm:r}
Let $G$ take values almost surely in an output subspace of real dimension at most $r$, with $\norm G\le B$ and input Hoeffding degree at most two. Set
$V(G)=\E\norm{G-\E G}^2$ and $\Inf_i(G)=\E\norm{D_iG}^2$.
For every product law,
\begin{equation}\label{eq:R}
\boxed{\displaystyle V(G)^2\le 3\Ctwo\,rB^2\max_i\Inf_i(G).}
\end{equation}
\end{theorem}

\begin{theorem}[Rank-optimal adaptive observation]\label{thm:obs}
Let $F$ be a known real-Hilbert-valued observable on a \emph{finite-coordinate} independent product, with $\norm F\le B$, input degree at most two, and $n$ coordinates. Let $\tau_r(F)$ denote the sum of covariance eigenvalues beyond the largest $r\ge1$ output directions. For each $\eta>0$ there is a deterministic adaptive policy of depth at most
\begin{equation}\label{eq:depth}
\boxed{\displaystyle N\le\min\left\{n,\left\lceil\frac{12\Ctwo rB^2}{\eta}\right\rceil\right\}}
\end{equation}
whose conditional-mean predictor has squared Bayes risk at most $\tau_r(F)+\eta$. Conversely, for every $r,k\ge1$ there is a $B=1$, rank-$r$, degree-one observable with $\tau_r=0$ such that achieving squared risk at most $\eta=1/(2k)$ requires at least $r/(4\eta)$ whole-symbol observations, even adaptively. The order $r/\eta$ in the universal upper bound is therefore optimal.
\end{theorem}

The self-contained degree-two proof and rank/adaptive consequences, which do not assume any all-degree extension, occupy Sections~\ref{sec:quad}--\ref{sec:policy}. The impossibility of several naive Boolean reductions is discussed in Section~\ref{sec:attack}; the distinct all-degree inequalities, their finite Lean premise, and the analytic lifting argument are presented only afterward in Section~\ref{sec:alldeg}. Neither a whole-symbol query-depth bound nor the quantum consequence is a polynomial-time implementation guarantee.

\section{Proof of the degree-two influence theorem}\label{sec:quad}
We prove Theorem~\ref{thm:q} on a finite full-support product. Bounded finite-coordinate/finite-partition conditional projections then give the general form in Section~\ref{sec:lift}. Write
\begin{equation}\label{eq:dec}
f=\E f+\sum_i h_i(X_i)+\sum_{i<j}h_{ij}(X_i,X_j),\quad
E_i h_i=0,\quad E_i h_{ij}=E_jh_{ij}=0.
\end{equation}
For each $i$, put $H_i=L^2(\mu_i;\mathbb R)$ and
\begin{equation}\label{eq:ui}
U_i(\cdot;X_{-i})=h_i(\cdot)+\sum_{j\ne i}h_{ij}(\cdot,X_j),\quad
\Inf_i(f)=\norm{h_i}_2^2+\sum_{j\ne i}\norm{h_{ij}}_2^2.
\end{equation}
For a pair $(i,j)$ set $b_{ij}=\max_{x_i,x_j}\abs{h_{ij}(x_i,x_j)}$.

\subsection{The disjoint-edge coupling budget}
\begin{lemma}[Matching mass]\label{lem:matching}
For any matching $\mathsf P$ in the graph of nonzero quadratic interactions,
$\sum_{(i,j)\in\mathsf P}b_{ij}\le4B$.
\end{lemma}
\begin{proof}
Fix one edge $(i,j)$ and a pair $(a,b)$ realizing $\abs{h_{ij}(a,b)}=b_{ij}$. Since $h_{ij}$ is centered separately in both inputs, the average of
$\Delta_{ij}=h_{ij}(a,b)-h_{ij}(a',b)-h_{ij}(a,b')+h_{ij}(a',b')$
over independent $a'\sim\mu_i$, $b'\sim\mu_j$ equals $h_{ij}(a,b)$. Choose $a',b'$ so that $\abs{\Delta_{ij}}\ge b_{ij}$.
For the chosen edge compare the \emph{even} two-point coupling, uniform on $(a,b),(a',b')$, to the \emph{odd} coupling, uniform on $(a,b'),(a',b)$. They share the same one-coordinate marginals and their $h_{ij}$ expectations differ by $\Delta_{ij}/2$. Independently make this comparison on every edge of the matching, orienting each difference positively; give all unmatched coordinates the same law in both experiments. All additive terms and all interactions between distinct blocks have identical expectations because the marginals agree and blocks are independent. Hence the difference of the two expectations of $f$ is $\frac12\sum_{e\in\mathsf P}\abs{\Delta_e}$. Both expectations lie in $[-B,B]$, so this difference is at most $2B$, proving the lemma.
\end{proof}

\subsection{Two norm budgets from pointwise boundedness}
\begin{lemma}[Additive oscillation and coordinate-section budgets]\label{lem:budgets}
For the decomposition \eqref{eq:dec},
\begin{equation}\label{eq:budgets}
\sum_i\norm{h_i}_2\le 2(1+\sqrt2)B,\qquad
\sum_i\E\norm{U_i(\cdot;X_{-i})}_{H_i}\le8B.
\end{equation}
\end{lemma}
\begin{proof}
For the first inequality consider the product noise operator $T_t=\prod_i(E_i+tD_i)$, $0\le t\le1$. It is a positive Markov contraction, and
$T_t f(x)=\E f+t\sum_i h_i(x_i)+t^2\sum_{i<j}h_{ij}(x_i,x_j)$ lies in $[-B,B]$.
For two inputs $x,y$, the difference $r(t)=T_tf(x)-T_tf(y)=ut+vt^2$ satisfies $r(0)=0$, $\abs{r(t)}\le2B$, and $u=\sum_i[h_i(x_i)-h_i(y_i)]$. At $s=\sqrt2-1$, exact interpolation gives
$u=[r(s)-s^2r(1)]/[s(1-s)]$, whence $\abs u\le2B(1+s^2)/[s(1-s)]=4(1+\sqrt2)B$.
The oscillation of the additive sum is the sum of its coordinate oscillations on the Cartesian support. Since every centered scalar random variable in an interval of width $w$ has standard deviation at most $w/2$, $\sum_i\norm{h_i}_2\le\frac12\osc(\sum h_i)\le2(1+\sqrt2)B$.

For the second inequality, fix $x$ and in each fiber choose a symbol maximizing $f(x_{-i},\cdot)-f(x)$, with gain $a_i^+\ge0$, and one minimizing it, with loss $a_i^-\ge0$. Independently switch each coordinate from its current symbol to its selected replacement with probability $t$. The resulting expected value $q_\pm(t)$ is a quadratic polynomial bounded by $B$ for $0\le t\le1$, and $q_+'(0)=\sum_i a_i^+$, $q_-'(0)=-\sum_i a_i^-$. The identity $q'(0)=-3q(0)+4q(1/2)-q(1)$ yields $\sum_i a_i^\pm\le8B$. Thus the sum of fiber oscillations at every $x$ is at most $16B$. Conditional on $X_{-i}$, the centered section $U_i$ is the fiber value minus its mean, so its $H_i$-norm is at most half its range. Integrating and summing proves the second inequality.
\end{proof}

\subsection{Greedy edge assignment and a fourth-moment inequality}
\begin{lemma}[Quadratic energy bound]\label{lem:energy}
Under \eqref{eq:dec}, with $M=\max_i\Inf_i(f)$,
\begin{equation}\label{eq:energy}
\sum_{i<j}\norm{h_{ij}}_2^2\le
\left[8+\sqrt3\bigl(8+2(1+\sqrt2)\bigr)\right]B\sqrt M.
\end{equation}
\end{lemma}
\begin{proof}
Process edges in nonincreasing order of $b_{ij}$, retaining an edge if disjoint from previously retained edges. The retained edges form a matching. Give both endpoints of a retained edge the cost $c_i=b_{ij}$ and give every unused vertex cost zero. Each original edge $(i,j)$ can be assigned to an endpoint with $c_i\ge b_{ij}$: if not retained, it was blocked by an earlier retained edge of at least that weight. The matching lemma implies $\sum_i c_i\le8B$.

Let $S_i$ consist of the neighbors of edges assigned to $i$ and put, in the Hilbert space $H_i$,
$Y_{ij}=h_{ij}(\cdot,X_j)$, $Z_i=\sum_{j\in S_i}Y_{ij}$, and
$s_i=\E\norm{Z_i}_{H_i}^2=\sum_{j\in S_i}\norm{h_{ij}}_2^2\le\Inf_i(f)\le M$.
For a fixed $i$ these $H_i$-valued random variables are independent and centered, with $\norm{Y_{ij}}_{H_i}\le b_{ij}\le c_i$. Expanding the fourth moment and using Cauchy--Schwarz,
\begin{equation}\label{eq:4moment}
\E\norm{Z_i}_{H_i}^4\le c_i^2s_i+3s_i^2.
\end{equation}
Indeed the diagonal terms satisfy $\E\norm{Y_{ij}}^4\le c_i^2\E\norm{Y_{ij}}^2$, while the pairwise inner-product terms are at most products of their second moments. The interpolation $\norm W_{L^2}^3\le\norm W_{L^1}\norm W_{L^4}^2$ applied to $W=\norm{Z_i}$ gives
\begin{equation}\label{eq:interpolate}
s_i\le\sqrt{c_i^2+3s_i}\,\E\norm{Z_i}_{H_i}
\le(c_i+\sqrt{3M})\E\norm{Z_i}_{H_i}.
\end{equation}
Since $Z_i$ is the conditional expectation of $U_i-h_i$ over the assigned neighbor coordinates, Jensen and the triangle inequality show
$\E\norm{Z_i}\le\E\norm{U_i}+\norm{h_i}_2$.
Therefore Lemma~\ref{lem:budgets} gives
\begin{equation}\label{eq:Zbudget}
\sum_i\E\norm{Z_i}\le\bigl[8+2(1+\sqrt2)\bigr]B,
\qquad \E\norm{Z_i}\le\sqrt{s_i}\le\sqrt M.
\end{equation}
Every edge is assigned exactly once, so $\sum_{i<j}\norm{h_{ij}}_2^2=\sum_i s_i$. Sum \eqref{eq:interpolate}, use $\sum c_i\le8B$ and \eqref{eq:Zbudget}, and obtain \eqref{eq:energy}.
\end{proof}

\begin{proof}[Proof of Theorem~\ref{thm:q} for finite full-support products]
The additive energy satisfies
\[\sum_i\norm{h_i}_2^2\le\sqrt M\sum_i\norm{h_i}_2\le2(1+\sqrt2)B\sqrt M.\]
By Hoeffding orthogonality and Lemma~\ref{lem:energy},
\[
\begin{aligned}
V&=\sum_i\norm{h_i}_2^2+\sum_{i<j}\norm{h_{ij}}_2^2\\
&\le\bigl[2(1+\sqrt2)+8+\sqrt3(8+2(1+\sqrt2))\bigr]B\sqrt M\\
&=D_2B\sqrt M.
\end{aligned}
\]
Squaring proves $M\ge V^2/(\Ctwo B^2)$.
\end{proof}

\begin{proof}[Proof of Proposition~\ref{prop:sharp}]
Let $Z_1,\ldots,Z_{2m}$ be independent fair signs and $f_m=m^{-1}\sum_{k=1}^m Z_{2k-1}Z_{2k}$. The products are independent fair signs, so $\abs{f_m}\le1$ and $\Var(f_m)=1/m$. Every coordinate occurs in one quadratic monomial of coefficient $1/m$, so its resampling influence is exactly $1/m^2$. A nonzero degree-two Fourier coefficient witnesses exact degree two.
\end{proof}

\section{Proof of the rank-linear degree-two theorem}\label{sec:rank}
The scalar proof is stable under Hilbert-valued interaction increments because it uses independent fourth moments rather than independent scalar coordinates of the output.

\begin{lemma}[Finite-sign Hilbert compression]\label{lem:sign}
For $w_1,\ldots,w_r$ in any real Hilbert space $\cK$ and independent fair signs $\varepsilon_j$,
\begin{equation}\label{eq:sign}
\E_\varepsilon\norm{\sum_{j=1}^r\varepsilon_jw_j}_{\cK}
\ge\frac{1}{\sqrt3}\left(\sum_{j=1}^r\norm{w_j}_{\cK}^2\right)^{1/2}.
\end{equation}
\end{lemma}
\begin{proof}
Let $W=\sum_j\varepsilon_jw_j$ and $s=\sum_j\norm{w_j}^2$. Independence gives $\E\norm W^2=s$ and
\[\begin{aligned}
\E\norm W^4&=\sum_j\norm{w_j}^4+2\sum_{j<k}\norm{w_j}^2\norm{w_k}^2\\
&\quad+4\sum_{j<k}\langle w_j,w_k\rangle^2\le3s^2.
\end{aligned}\]
Interpolation yields $s^{3/2}\le(\E\norm W)(\E\norm W^4)^{1/2}\le\sqrt3\,s\E\norm W$. Divide by $s$, with $s=0$ trivial.
\end{proof}

\begin{proof}[Proof of Theorem~\ref{thm:r} on finite full-support products]
Identify the effective output subspace with $\mathbb R^r$ and decompose $G$ into vector-valued $h_i,h_{ij}$. For every sign vector $\varepsilon\in\{-1,1\}^r$ the scalar $g_\varepsilon=\langle\varepsilon,G\rangle/\sqrt r$ satisfies $\abs{g_\varepsilon}\le B$. Apply Lemmas~\ref{lem:matching} and \ref{lem:budgets} to every such scalar function and average in $\varepsilon$.

For the matching bound, fix one point attaining $b_{ij}=\max\norm{h_{ij}}$ on each edge of a prescribed matching before choosing $\varepsilon$. Lemma~\ref{lem:sign}, first in $\mathbb R^r$, gives
\begin{equation}\label{eq:vectorbudgets}
\sum_{(i,j)\in\mathsf P}b_{ij}\le4B\sqrt{3r}.
\end{equation}
For the additive and coordinate-section budgets apply the same sign lemma in the spaces $L^2(\mu_i)$ (to the $r$ scalar components of $h_i$) and $L^2(\mu_i)$ (to each fiber of $U_i$), then integrate. The result is
\begin{equation}\label{eq:vectorbudget2}
\sum_i\norm{h_i}_2\le2(1+\sqrt2)B\sqrt{3r},\qquad
\sum_i\E\norm{U_i(\cdot;X_{-i})}_{L^2(\mu_i;\mathbb R^r)}\le8B\sqrt{3r}.
\end{equation}
Run the identical descending-matching assignment with vector-valued $b_{ij}$. It yields $\sum_i c_i\le8B\sqrt{3r}$. The assigned increments $Y_{ij}=h_{ij}(\cdot,X_j)$ are independent centered random elements of $L^2(\mu_i;\mathbb R^r)$, with $\norm{Y_{ij}}\le c_i$. Thus the exact same fourth-moment and interpolation argument \eqref{eq:4moment}--\eqref{eq:interpolate} applies, and
\[
\sum_{i<j}\norm{h_{ij}}_2^2\le
\left[8+\sqrt3\bigl(8+2(1+\sqrt2)\bigr)\right]B\sqrt{3rM(G)}.
\]
The additive energy is at most $2(1+\sqrt2)B\sqrt{3rM(G)}$. Orthogonality gives $V(G)\le D_2 B\sqrt{3rM(G)}$, hence \eqref{eq:R}.
\end{proof}

\begin{remark}[Why the rank factor is structurally meaningful]
The argument does not assert that the individual components of $G$ are independent. The $\sqrt r$ loss comes only from compressing bounded Hilbert outputs to bounded random scalar projections, while the assigned neighbor increments are independent because their \emph{input coordinates} are independent. The lower bound in Section~\ref{sec:policy} shows that the induced linear $r$ dependence in observation depth cannot generally be eliminated, although the constant is not claimed optimal.
\end{remark}

\section{Proof of rank-optimal adaptive quadratic observation}\label{sec:policy}
The key fact for adaptive observation is not merely that each posterior has an influential coordinate. It is that a \emph{degree-weighted residual energy} drops by exactly that coordinate's influence in conditional expectation.

\begin{lemma}[Actual degree-weighted posterior drift]\label{lem:drift}
Let $G$ be Hilbert-valued of input degree at most two on a finite product. At a transcript revealing coordinates $Q$, decompose the remaining-input section into orthogonal Hoeffding components $G_S$ and write
\[
v_Q=\sum_{S\ne\varnothing}\norm{G_S}_2^2,\qquad
p_Q=\sum_{S\ne\varnothing}|S|\norm{G_S}_2^2=\sum_{i\notin Q}\Inf_i(G\mid Q).
\]
Then $v_Q\le p_Q\le2v_Q$. If an unobserved coordinate $i$ is revealed, conditional on the current transcript,
\begin{equation}\label{eq:drift}
p_Q-\E[p_{Q\cup\{i\}}\mid\text{current transcript}]=\Inf_i(G\mid Q).
\end{equation}
\end{lemma}
\begin{proof}
For each nonempty residual support $S$ not containing $i$, the contribution $G_S$ persists. A term $G_{S\cup\{i\}}$ that contains $i$ becomes, after observation of $X_i$, a random coefficient on $S$, decreasing its degree weight by one. Terms now sharing a residual support $S$ have cross-inner-product expectation zero over the newly revealed $X_i$, because $E_i G_{S\cup\{i\}}=0$. Hence the expected new weighted energy is
$\sum_{S:\,i\notin S}|S|\norm{G_S}_2^2 +\sum_{S:\,i\in S}(|S|-1)\norm{G_S}_2^2$,
where the term with $S=\{i\}$ has new weight zero. Subtract from the old energy to get $\sum_{S\ni i}\norm{G_S}_2^2=\Inf_i(G\mid Q)$. The sandwich follows because each positive-degree term has degree one or two.
\end{proof}

\begin{proof}[Proof of the upper bound in Theorem~\ref{thm:obs}]
Let $P_r$ project onto $r$ leading covariance eigenvectors of $F$ and set $G=P_rF$, so $\norm G\le B$, degree at most two, and effective output dimension at most $r$. At every possible transcript the unrevealed coordinates retain their original independent laws and the section of $G$ is still bounded and quadratic. Theorem~\ref{thm:r} therefore gives, with $K_r=3\Ctwo rB^2$,
\[
m_Q:=\max_{i\notin Q}\Inf_i(G\mid Q)\ge\frac{v_Q^2}{K_r}\ge\frac{p_Q^2}{4K_r}.
\]
Choose at each transcript a coordinate attaining $m_Q$, breaking ties in a fixed way. Freeze branches after all coordinates are queried. Let $P_t=\E[p_Q]$ at observation depth $t$. By Lemma~\ref{lem:drift} and Jensen,
\begin{equation}\label{eq:recurrence}
P_t-P_{t+1}=\E[m_Q]\ge\frac{\E[p_Q^2]}{4K_r}\ge\frac{P_t^2}{4K_r}.
\end{equation}
Until $P_t=0$, reciprocal telescoping gives $P_t\le4K_r/t$ for $t\ge1$. The conditional-mean predictor of $G$ has squared risk $\E v_Q\le P_t$. Thus $t\ge4K_r/\eta=12\Ctwo rB^2/\eta$ suffices. Reading all $n$ coordinates gives zero error for $G$, producing the cap in \eqref{eq:depth}.

The component $F-G$ is pointwise orthogonal to $G$ in the fixed output decomposition. Its posterior conditional variance is at most its unconditional variance $\tau_r(F)$ by the law of total variance, and the cross term in conditional squared norm vanishes by the output-space orthogonality. The total risk is at most $\tau_r(F)+\eta$.
\end{proof}

\begin{proof}[Proof of the lower bound in Theorem~\ref{thm:obs}]
Let $Z_{jt}$ ($1\le j\le r$, $1\le t\le k$) be independent fair signs and define
\begin{equation}\label{eq:lower}
F(Z)=\frac1{k\sqrt r}\left(\sum_{t=1}^kZ_{1t},\ldots,\sum_{t=1}^kZ_{rt}\right)\in\mathbb R^r.
\end{equation}
Every input satisfies $\norm F\le1$, the function has degree one, its covariance rank is $r$ and its total variance is $1/k$. Each observed sign reduces the conditional squared variance by exactly $1/(rk^2)$, regardless of other observations or adaptivity. To reduce the risk to $1/(2k)$ requires at least $rk/2=r/(4\eta)$ observations.
\end{proof}

\section{Why direct Boolean reductions do not settle the product-law question}\label{sec:attack}
\subsection{A support-degree obstruction to deterministic fair-bit coding}
\begin{lemma}[Support of a low-degree Boolean polynomial]\label{lem:support}
A nonzero multilinear polynomial $h:\{0,1\}^m\to\mathbb R$ of degree at most $s$ is nonzero at at least $2^{m-s}$ inputs.
\end{lemma}
\begin{proof}
Induct on $m$. If $h$ does not depend on the last bit, double the inductive support count. Otherwise write $h(x,t)=A(x)+tB(x)$ with $B\not\equiv0$ and $\deg B\le s-1$. Wherever $B(x)\ne0$, at least one of $A(x)$ and $A(x)+B(x)$ is nonzero. The inductive count for $B$ is $2^{(m-1)-(s-1)}=2^{m-s}$.
\end{proof}
If a deterministic map of $m$ fair bits produces a designated symbol with probability $p>0$, its indicator has support $p2^m$ and degree at least $\log_2(1/p)$. A rare Bernoulli symbol of probability $2^{-m}$ therefore costs degree $m$ under every exact deterministic fair-bit coding (the AND function attains this). A symbol-level degree-$d$ observable need not retain degree $d$ after encoding: straightforward composition may have degree $dm$ and a cube $d^{-6}$ bound becomes $(dm)^{-6}$. This is an obstruction to that \emph{coding route}, not a theorem prohibiting all reductions.

\subsection{The maximum cannot be moved through random-subcube averaging}
Draw independent copies $X,Y$ of the product input, condition on their unordered coordinate pairs, and independently choose each orientation by a fair sign. The resulting Boolean-cube restriction $g_\Pi$ has degree at most $d$, and, writing $J_i(\Pi)$ for its conditional fair-bit influence,
\begin{equation}\label{eq:subcube}
\E_\Pi J_i(\Pi)=\tfrac12\Inf_i(f),\qquad
\E_\Pi\Var(g_\Pi)=\sum_{S\ne\varnothing}(1-2^{-\abs S})\norm{f_S}_2^2\ge\tfrac12\Var(f).
\end{equation}
Turturean's theorem on each subcube bounds $\E\max_i J_i$, not $\max_i\E J_i$. The distinction cannot be suppressed by a universal constant.

\begin{proposition}[Exact maximum/expectation obstruction]\label{prop:exchange}
Let $n\ge2$, let $X_i\sim\operatorname{Bernoulli}(1/n)$ independently and $f(X)=n^{-1}\sum_i X_i$. Then $f$ is bounded and has degree one, yet
\begin{equation}\label{eq:exchange}
\frac{\E_\Pi\max_i J_i(\Pi)}{\max_i\Inf_i(f)}
=\frac{1-\left(1-2\frac1n(1-\frac1n)\right)^n}{4\frac1n(1-\frac1n)}
\sim\frac{1-e^{-2}}4\,n.
\end{equation}
\end{proposition}
\begin{proof}
The original influence is $p(1-p)/n^2$, where $p=1/n$. The conditional fair-bit influence $J_i$ is $1/(4n^2)$ exactly when $X_i\ne Y_i$, and zero otherwise. These events are independent with probability $2p(1-p)$. Thus the expected maximum is $[1-(1-2p(1-p))^n]/(4n^2)$, giving \eqref{eq:exchange}.
\end{proof}
The constructive lesson is limited but important: a product transfer must preserve a \emph{common weighted inequality before averaging}, instead of selecting a maximizing coordinate separately within each random subcube. The all-degree companion of Section~\ref{sec:alldeg} uses precisely this order. In the degree-two argument below we avoid this issue altogether by using the actual interaction graph rather than conditioning on random subcubes.


\section{One-query quantum observation and the adaptivity boundary}\label{sec:quantum}
Consider the usual coherent \emph{whole-symbol} oracle on $n$ finite-alphabet coordinates and one quantum query. Every output amplitude is a linear combination of functions reading at most one coordinate. Squaring amplitudes and summing measurement probabilities yields observables of interaction degree at most two; this is the whole-symbol version of the standard polynomial method \cite{BBCMW2001}.

\begin{corollary}[Prediction of one-query quantum outcomes]\label{cor:quantum}
If $P(X)=(P_1(X),\ldots,P_s(X))$ is the output probability vector of such an algorithm, with independent input symbols under an arbitrary known product law, then for every $\eta>0$ there is a deterministic adaptive classical policy that reads at most
\[
\min\left\{n,\left\lceil\frac{12\Ctwo s}{\eta}\right\rceil\right\}
\]
whole input symbols and returns a probability-vector estimate $\widehat P$ satisfying $\E\norm{P(X)-\widehat P(X)}_2^2\le\eta$. No restriction is made on the policy's computation time or representation size.
\end{corollary}
\begin{proof}
Each $P_j$ has interaction degree at most two and $\norm{P(x)}_2\le\sum_jP_j(x)=1$. Apply Theorem~\ref{thm:obs} with $F=P$, $B=1$, $r=s$, and $\tau_s=0$. Posterior conditional expectations of probability vectors remain probability vectors.
\end{proof}

\begin{corollary}[Scalar one-query prediction with an error probability]\label{cor:scalarq}
If $g(X)\in[0,1]$ is a one-query acceptance probability, then for each $0<\varepsilon,\delta<1$ there is a deterministic adaptive predictor $\widehat g$ with
\[
\Pr\{\abs{g(X)-\widehat g(X)}>\varepsilon\}\le\delta,
\qquad N\le\min\left\{n,\left\lceil\frac{\Ctwo}{\delta\varepsilon^2}\right\rceil\right\}.
\]
\end{corollary}
\begin{proof}
Apply the scalar case of Theorem~\ref{thm:obs} to $f=2g-1$, obtaining $\E\abs{f-\widehat f}^2\le4\delta\varepsilon^2$ after at most $\lceil4\Ctwo/(4\delta\varepsilon^2)\rceil$ observations. Use $\widehat g=(\widehat f+1)/2$ and Markov's inequality.
\end{proof}

\begin{remark}[Adaptivity is necessary for a uniform two-query phenomenon]
Let $A$ be uniform on $[m]$ and $Z_1,\ldots,Z_m$ be independent fair signs; set $f=Z_A$. The function is bounded and has degree two. Querying $A$ and then $Z_A$ recovers it exactly in two adaptive queries. Querying $A$ and $q$ \emph{predetermined} leaves leaves mean-squared error $1-q/m$; querying $q$ leaves but not $A$ leaves error $1-q/m^2$. Thus the dimension- and alphabet-independent adaptive conclusion should not be advertised as a corresponding nonadaptive theorem. No originality claim is made for this standard addressing example.
\end{remark}

\section{From finite alphabets to arbitrary products}\label{sec:lift}
The extension is a contraction argument, not the new analytic heart. We include it to make the exact scope of the theorems unambiguous.

\begin{lemma}[Bounded product approximation and influence contraction]\label{lem:lift}
Let $J$ be any index set and $f\in L^2(\bigotimes_{i\in J}\mu_i)$ be bounded by $B$ with Hoeffding interaction degree at most $d$. Finite-coordinate conditional expectations, followed by conditional expectations onto finite coordinatewise measurable partitions, are bounded by $B$, preserve degree at most $d$, converge to $f$ in $L^2$, and do not increase any resampling influence. For every finite $F\subset J$, $\sum_{i\in F}\Inf_i(f)\le d\Var(f)$.
\end{lemma}
\begin{proof}
Finite rectangles generate the product sigma-algebra and their linear span is dense in product $L^2$, even for an arbitrary index set, using the directed system of finite-cylinder subspaces. Every bounded conditional expectation is an $L^2$ contraction and preserves its essential bound. Conditional integration in coordinate $i$ commutes with conditional expectation onto finite products of coordinate partitions, so $D_i P=PD_i$ and $\Inf_i(Pf)=\norm{PD_if}_2^2\le\Inf_i(f)$. By the Hoeffding representation, $P$ preserves the linear span and closure of interactions on at most $d$ coordinates. Hence variance and all fixed-coordinate influences converge along the directed net. The last inequality follows from \eqref{eq:anova} for finite projected products and passage to the limit.
\end{proof}
\begin{proof}[Completion of Theorems~\ref{thm:q} and \ref{thm:r}]
Apply the finite-alphabet inequalities to the bounded finite-partition projections and let their variances tend to $V$. Influence contraction yields the same inequality with $M=\sup_i\Inf_i(f)$, with no constant loss. When $V>0$, this supremum is positive. The superlevel set $\{i:\Inf_i(f)>M/2\}$ is finite: otherwise finite subsets of arbitrary size contradict $\sum_{i\in F}\Inf_i(f)\le dV$. Its largest influence must equal $M$, so the supremum is attained. For Hilbert-valued functions the same proof applies to vector $L^2$ and a fixed finite output subspace. Completed product sigma-algebras do not change the argument after replacing a bounded function by an almost-surely equal product-measurable representative.
\end{proof}

\section{Complementary all-degree influence: separate statement and proof dependencies}\label{sec:alldeg}
The all-degree and quadratic theorems are \emph{different conclusions}. The all-degree exponent $5/2$ and $d^{-6}$ dependency extend Turturean's cube framework to arbitrary independent source laws. The degree-two matching argument yields a stronger variance exponent but does not establish a degree-$d$ exponent-two inequality for $d>2$. We state precisely what the recorded finite Lean theorem proves, give the law-uniform lifting argument for its weaker coefficient, and isolate the extra analytic work behind the sharper v8 coefficient.

\begin{theorem}[All-degree bias-free product influence; sharpened analytic coefficient]\label{thm:alldeg}
Let $\{X_i\}_{i\in J}$ be any independent family of random elements on arbitrary probability spaces, with $J$ of any cardinality. Let $f\in L^2(\mu)$ be real, $|f|\le B$ almost surely, $B>0$, with Hoeffding interaction degree at most $d\ge1$ and $V=\Var(f)>0$. The v8 analytic product-law argument gives an attained maximal influence satisfying
\begin{equation}\label{eq:alldegheadline}
\boxed{\displaystyle \max_{i\in J}\Inf_i(f)\ge
 \cprod\frac{V^{5/2}}{B^3d^6},\qquad
 \cprod=\frac{(8/5)^{3/2}\sqrt{15}}{125(4-2\sqrt3)\pi^4}
 =0.001201250694\ldots.}
\end{equation}
If every effective nontrivial coordinate alphabet is binary, the source gives the larger coefficient
$2(8/5)^{3/2}\sqrt{15}/(125\pi^4)=0.001287496613\ldots$.
The coefficients have no bias, minimum-atom, alphabet-size or coordinate-count dependence. \emph{Evidence status:} detailed author-side analytic proof in the cited v8 source; not an end-to-end Lean certificate or independent referee determination.
\end{theorem}

\begin{theorem}[All-degree bias-free influence; certified finite premise and analytic lift]\label{thm:lean}
Under the same assumptions, a separately grounded all-degree inequality is
\begin{equation}\label{eq:lean}
\boxed{\displaystyle \max_{i\in J}\Inf_i(f)\ge c_*\frac{V^{5/2}}{B^3d^6},\qquad
 c_*:=\frac{\sqrt{3/12500}}{\pi^6}=0.0000161141122537\ldots.}
\end{equation}
For a finite product represented on a common finite alphabet with nonnegative normalized marginal tables and $|f|\le B$ everywhere, the normalized inequality is the consequence of the named Lean declaration 
\texttt{Nightfall.ProductSpace.Influence.product\_influence}, whose hypotheses include bounded indicator-monomial degree $\le d$. The arbitrary-product statement follows by a separate analytic finite-partition argument proved in Section~\ref{sec:alldeg}. Neither the continuum lift nor the larger coefficient in Theorem~\ref{thm:alldeg} was compiled by that declaration.
\end{theorem}


\begin{lemma}[Exact finite statement of the pinned kernel theorem]\label{lem:finitelean}
Fix finite index and alphabet types $I,A$, nonnegative tables $p_i(a)$ with $\sum_a p_i(a)=1$, and $g:A^I\to[-1,1]$ of bounded indicator-monomial degree at most $d\ge1$. Write $V=\Var_p(g)>0$ and define resampling influences by~\eqref{eq:infdef}. Then the formally recorded declaration proves
\begin{equation}\label{eq:finitelean}
 \exists j\in I:\quad
 3V^5\le12500(\pi^2d^2)^6\Inf_j(g)^2,
 \qquad\text{hence}\qquad
 \Inf_j(g)\ge c_*\frac{V^{5/2}}{d^6}.
\end{equation}
Marginal tables may assign zero mass to symbols; the degree is that of the explicit finite indicator representation. The certified premise is \emph{not} a statement about arbitrary measurable input spaces.
\end{lemma}
\begin{proof}
The first inequality is the source- and commit-pinned Lean declaration quoted in Section~\ref{sec:status} (and is imported here, not re-proved by the present manuscript). Since influences and variances are nonnegative, take square roots and divide by $\sqrt{12500}\,\pi^6d^6$. This yields $\sqrt{3/12500}\,V^{5/2}/(\pi^6d^6)$ exactly.
\end{proof}

\begin{proof}[Analytic proof of the arbitrary-product lift in Theorem~\ref{thm:lean}]
First normalize $h=f/B$, so $|h|\le1$, $\Var(h)=V/B^2$, and $\Inf_i(h)=\Inf_i(f)/B^2$. Let $P_\alpha$ be the net of coordinatewise finite-partition, finite-coordinate conditional expectations from Lemma~\ref{lem:lift}. Then $g_\alpha=P_\alpha h$ is finite-valued, bounded by one, and supported on finitely many coordinates. Conditional averaging preserves the interaction-degree cap and each resampling influence contracts:
\[
 \Var(g_\alpha)\longrightarrow\Var(h),\qquad
 \Inf_i(g_\alpha)\le\Inf_i(h)\quad(i\in J).
\]
On each finite partition, represent $g_\alpha$ using the indicators of cells on at most $d$ coordinate sets. This representation has indicator-monomial degree at most $d$ because its Hoeffding expansion has no interaction above order $d$. Distinct finite alphabets can be embedded into a common finite alphabet by zero-mass padding. Define the padded function by mapping each unavailable symbol to a fixed valid symbol in its coordinate, retaining the pointwise $[-1,1]$ bound, interaction degree, and values on positive-mass inputs. Thus Lemma~\ref{lem:finitelean} applies with the same $d$ and no minimum-atom loss.

Let $M=\sup_{i\in J}\Inf_i(h)$. Each finite instance therefore obeys
\[
 M\ge\max_i\Inf_i(g_\alpha)
 \ge c_*\frac{\Var(g_\alpha)^{5/2}}{d^6}.
\]
Taking the directed limit yields $M\ge c_*\Var(h)^{5/2}/d^6>0$. Moreover, for every finite $F\subset J$, Lemma~\ref{lem:lift} gives $\sum_{i\in F}\Inf_i(h)\le d\Var(h)$. Therefore the superlevel set $\{i:\Inf_i(h)\ge M/2\}$ is finite and nonempty; a maximum on this set equals the full supremum. Rescaling back by $B^2$ gives exactly~\eqref{eq:lean}. Completion of the product sigma-algebra does not alter the law of any almost-sure product-measurable representative or its resampling influences.
\end{proof}

\paragraph{Stronger all-degree constant: the additional analytic operator chain.}
Theorem~\ref{thm:alldeg} does \emph{not} follow from the weaker finite Lean theorem merely by rescaling or approximating. Its source proof in v8~\cite{WrightV8} separately derives these nontrivial finite results: (i) a weighted all-pairs reciprocal inequality by conditioning two independent product draws on unordered pairs, applying Turturean's antipodal theorem with weights fixed \emph{before} conditioning; (ii) localized Cauchy--Schwarz weight removal and a two-sided level-set/median moment estimate; (iii) a competing-symbol Chebyshev selector with an LP-dual row cost at most $\pi^2d^2/2$ independent of alphabet size; and (iv) sharp iid fibre resampling with $\rho=4-2\sqrt3$, giving an aggregate increment estimate proportional to $4\rho d^2$. The sharpened constant is obtained only after composing this chain and lifting it with finite partitions. The v8 source supplies the detailed arguments; the present v9.2 is not a substitute for independent verification of their missing proof gates.

\begin{remark}[What the reduction audit establishes]
Girish--Mittal--Nehoran--Raz's fixed-bias gadget \cite[Theorem~4.3 and Remarks~4.4--4.5]{GMNR2026} incurs bias-dependent overhead; exact fair-bit coding can inflate degree; and Proposition~\ref{prop:exchange} excludes a universal exchange of $\E\max$ with $\max\E$. These facts identify failures of three straightforward proof routes, not a proof of global priority or the impossibility of a different reduction. The exact coefficient in the sharp resampling lemma is a distinct formal achievement and does not by itself certify the stronger all-degree theorem.
\end{remark}

\section{Scope boundaries, comparisons, and a fixed-set guarantee}\label{sec:comparisons}
\paragraph{The binary quadratic case already has stronger classical constants.}
For globally bounded quadratic $f:\{-1,1\}^n\to[-B,B]$, a Grothendieck duality argument (included in v8) gives $\sum_i\sqrt{\Inf_i(f)}\le4K_G^{\mathbb R}B$ under any product law, hence $\max_i\Inf_i(f)\ge V^2/[16(K_G^{\mathbb R})^2B^2]$. This is stronger than \eqref{eq:Q}; we do not claim a novel binary quadratic coefficient. The claim warranting review is the \emph{categorical/nonatomic} statement derived only from $\abs f\le B$, where neither a common two-point alphabet nor a completely bounded norm is assumed. The general Efron--Stein decomposition itself is classical \cite{Blais2008}; it is the dimension-free \emph{inequality} that is under consideration.

\paragraph{Completely bounded methods concern a different norm regime.}
Escudero Guti\'errez \cite{EG2023} proved variance-squared influence statements for homogeneous Fourier-completely-bounded polynomials, while Escudero Guti\'errez--Saucedo--Palazuelos \cite{EGSP2026} give sharp root-influence bounds for block-multilinear completely bounded forms. Those results should not be presented as nonexistent; the restricted norm hypothesis must be accounted for in any scope comparison. A dedicated check against general categorical bilinear norm theorems remains mandatory before originality is asserted.

\paragraph{Product independence cannot simply be removed.}
For $0<\epsilon<1$, let $\nu_\epsilon=(1-\epsilon)\Law(Z,Z)+\epsilon\Law(Z_1,Z_2)$ with independent fair signs $Z,Z_1,Z_2$, and $f(x_1,x_2)=x_1$. Each of four inputs has positive mass and the marginals are fair, but $\Var_{\nu_\epsilon}(f)=1$, $\Inf^{\nu_\epsilon}_2(f)=0$ and
$\Inf^{\nu_\epsilon}_1(f)=1-(1-\epsilon)^2=2\epsilon-\epsilon^2\to0$. Thus there is no arbitrary-correlated-law version with a positive variance-only influence floor. If a correlated $\nu$ has a reference product law $\mu$ with $0<a\le d\nu/d\mu\le b$, then the variational characterizations of conditional variance give
\begin{equation}\label{eq:correlated}
\Inf_i^\nu(f)\ge a\Inf_i^\mu(f),\qquad\Var_\nu(f)\le b\Var_\mu(f),
\quad\Longrightarrow\quad \max_i\Inf_i^\nu(f)\ge\frac{a}{b^2}\frac{\Var_\nu(f)^2}{\Ctwo B^2}
\end{equation}
for functions of product-reference degree at most two. This is a change-of-measure corollary, not a theorem with constants independent of arbitrary correlations.

\begin{proposition}[A fixed-set guarantee above the quadratic energy floor]\label{prop:nonadaptive}
For scalar $f$ of degree at most two with $\abs f\le B$ on a finite product, let $\delta_2=\sum_{|S|=2}\norm{f_S}_2^2$. For every $\eta>\delta_2$, a predetermined set of at most
\begin{equation}\label{eq:fixed}
\min\left\{n,\left\lceil\frac{(3+2\sqrt2)B^2}{\eta-\delta_2}\right\rceil\right\}
\end{equation}
coordinates yields conditional-mean risk at most $\eta$.
\end{proposition}
\begin{proof}
Let $b_i=\norm{h_i}_2$ for the additive Hoeffding terms. Lemma~\ref{lem:budgets} gives $L=\sum_i b_i\le2(1+\sqrt2)B$. Query the $q$ largest $b_i$. If $t=b_{q+1}$, then $\sum_{i>q}b_i^2\le t\sum_{i>q}b_i\le t(L-qt)\le L^2/(4q)\le(3+2\sqrt2)B^2/q$. All unrevealed quadratic energy is at most $\delta_2$, so the residual risk is at most $\delta_2+(3+2\sqrt2)B^2/q$. Take the stated ceiling or all coordinates, as appropriate.
\end{proof}

\section{Certification ledger, limitations, and priority}\label{sec:status}
\begin{center}\small
\begin{tabular}{@{}>{\raggedright\arraybackslash}p{0.32\linewidth}>{\raggedright\arraybackslash}p{0.35\linewidth}>{\raggedright\arraybackslash}p{0.24\linewidth}@{}}\toprule
\textbf{Claim} & \textbf{Evidence carried here} & \textbf{Certification state}\\\midrule
Quadratic scalar exponent two, all alphabets & Matching, oscillation, fourth-moment proof, Sections~\ref{sec:quad}, \ref{sec:lift} & Analytic proof; external audit pending\\
Hilbert rank-$r$ inequality & Sign-compression proof, Section~\ref{sec:rank} & Analytic proof; not Lean-certified\\
Adaptive rank-optimal depth & Exact posterior drift; $\Omega(r/\eta)$ example, Section~\ref{sec:policy} & Analytic proof; not Lean-certified\\
One-query quantum output & Polynomial method plus observation theorem & Conditional corollary of analytic theorem\\
All-degree stronger coefficient $\cprod$ & V8 weighted transfer, rows, median and finite lift & Analytic source; full-chain Lean and external review open\\
All-degree weaker $c_*$, arbitrary products & Lean finite premise plus explicit Section~\ref{sec:alldeg} $L^2$ lift & Finite premise kernel-checked; measurable lift analytic\\
Deterministic-encoding and max/expectation obstructions & Elementary proofs; exact script in reduction audit & Reproducible limited checks, not priority clearance\\\bottomrule
\end{tabular}
\end{center}

\noindent\textbf{New finite-instance audit.} An independently implemented exact-rational script accompanying this v9 recomputed the Hoeffding components directly from conditional tables on 24 heterogeneous $3$- and $4$-coordinate products with strictly positive $2$- and $3$-symbol marginals. It checked every matching budget, the scalar additive and fiber norm budgets with rigorous rational square-root enclosures, each assigned Hilbert fourth-moment inequality, the final quadratic inequality, and the actual posterior-energy drift after each possible first coordinate observation. All checks passed. This is a reproducible finite-instance attack, not a general theorem certificate or a substitute for independent proof review.\par\medskip

\noindent\textbf{Frozen certificate record (reported in v8, not rerun here).}
\begin{quote}\small\raggedright
\textbf{Declaration:} \texttt{Nightfall.ProductSpace.Influence.product\_influence}\\
\textbf{Pinned commit:} \texttt{\seqsplit{e67990dfc252b6ef9446c7b7c1a9ebf06e5c58eb}}\\
\textbf{Source blob:} \texttt{\seqsplit{70f69d8caebfaef40efdf16bab63f223753c190f}}\\
\textbf{Finite-product CI:} \url{https://github.com/d3lefts/d3lefts/actions/runs/42929857470}\\
\textbf{Heterogeneous-alphabet companion:} merged PR \#325, CI
\url{https://github.com/d3lefts/d3lefts/actions/runs/42932085109}.
\end{quote}
These identifiers reproduce the earlier source provenance. The present v9.2 does not claim to recompile or strengthen the frozen Lean certificate. The finite-to-general argument in Section~\ref{sec:alldeg} is separate from that receipt.\par\medskip

\noindent\textbf{What has not been proved here.} We do not establish the unrestricted degree-$d$ exponent-two conjecture $\max_i\Inf_i(f)\gtrsim V^2/d^4$, an efficient algorithm for computing the adaptive policy, a nonadaptive guarantee below the quadratic energy floor, an arbitrary-correlated-law theorem without comparability assumptions, or worldwide originality. The quadratic proof relies on genuinely independent increments after assigning pair-interaction edges. Extending its fourth-moment step to cubic interactions with overlapping cores requires a different mechanism; pairwise orthogonality alone does not imply the independent fourth-moment bound.

\noindent\textbf{Independent criticism gate.} Before circulation as a refereed contribution, a reader not involved in construction should (i) reproduce the even/odd coupling with complex alphabets and zero-mass limits, (ii) verify the fourth-moment coefficient and greedy assignment for arbitrary vector-valued kernels, (iii) audit measurable policy selection on non-finite alphabets, (iv) check the exact posterior energy formula on nontrivial mixed interactions, and (v) search scope-matched categorical bilinear norm literature. If a shorter categorical Grothendieck argument exists, the originality and headline must be revised accordingly.

\section{Conclusion}
This degree-two-first revision places the law-uniform quadratic influence theorem, its sharp exponent, its Hilbert rank bound and its rank-optimal adaptive consequence ahead of every general-degree claim. The all-degree extension is preserved later as a distinct theorem with separately documented proof dependencies. The stronger all-degree coefficient comes from the v8 weighted two-draw/row/median argument and remains subject to independent proof review; a weaker all-degree constant has a pinned finite Lean foundation and the explicit analytic lift given here. The quadratic matching, oscillation, and independent fourth-moment mechanism yields the sharp variance exponent at degree two, and its rank-sensitive posterior argument implies matching-order classical adaptive observation. These claims form a structured hierarchy with distinct premises, rather than one undifferentiated certification.

\begin{thebibliography}{99}\setlength{\itemsep}{3pt}
\bibitem{Turturean2026} D. Turturean, \emph{A proof of the Aaronson--Ambainis conjecture}, arXiv:2611.09773 (6 Nov 2026). \url{https://arxiv.org/html/2611.09773}.
\bibitem{GMNR2026} U. Girish, K. Mittal, B. Nehoran, and R. Raz, \emph{Distributional Variants of the Aaronson--Ambainis Conjecture}, arXiv:2609.35327 (28 Sep 2026). \url{https://arxiv.org/html/2609.35327}.
\bibitem{AA2014} S. Aaronson and A. Ambainis, \emph{The need for structure in quantum speedups}, Theory of Computing \textbf{10} (2014), 133--166.
\bibitem{BBCMW2001} R. Beals, H. Buhrman, R. Cleve, M. Mosca, and R. de Wolf, \emph{Quantum lower bounds by polynomials}, J. ACM \textbf{48} (2001), 778--797.
\bibitem{Blais2008} E. Blais, R. O'Donnell, and K. Wimmer, \emph{Polynomial regression under arbitrary product distributions}, COLT (2008), 193--204.
\bibitem{EG2023} F. Escudero Guti\'errez, \emph{Influences of Fourier completely bounded polynomials and classical simulation of quantum algorithms}, arXiv:2304.06713 (2023).
\bibitem{EGSP2026} F. Escudero Guti\'errez, M. Saucedo, and C. Palazuelos, \emph{Optimal inequalities for completely bounded polynomials and the limitations of quantum query algorithms}, arXiv:2609.05201 (4 Sep 2026).
\bibitem{WrightV8} D. J. Wright, \emph{Bias-Free Aaronson--Ambainis on All Product Spaces: Quadratic Influence and Rank-Optimal Observation}, NIGHTFALL specialist-referee draft v8, 11 Nov 2026 (32 pp., internal source manuscript and proof-passport ledger).
\bibitem{WrightAttack} NIGHTFALL, \emph{v8 adversarial reduction attack and surviving claims}, internal proof-level reduction report and exact-rational checker, 11 Nov 2026.
\end{thebibliography}
\end{document}
